\chapter{Conclusion}
\label{chap:conclusion}

This thesis has presented \gls{SafeVision}, a system for detecting and redacting sensitive information in visual media. The aim of the project was to design and develop a system that could automatically identify privacy-relevant content while supporting user review and control of the \gls{anonymization} process. The project was grounded in a legal and privacy oriented understanding of sensitive information based on GDPR, and this understanding was further informed by user testing and stakeholder feedback. This provided the basis for developing a practical solution for detecting and anonymizing sensitive visual content.

\gls{SafeVision} was completed as an interactive web application that combines automated detection, user-guided review, and redaction. The final solution consists of a React and TypeScript frontend, a FastAPI backend for communication, and Python-based computer vision components, including multiple \acrshort{ml} models and video tracking functionality. The system supports automatic detection of faces, text, license plates, barcodes, and selected object categories, while allowing users to review and adjust detections before final redaction is applied. It also provides multiple redaction options, giving users flexibility in how sensitive information is anonymized.

In conclusion, \gls{SafeVision} fulfills the project aim by providing a working system and functional foundation for privacy-preserving processing of visual media. Through its support for image and video processing, multiple detection categories, and flexible redaction options, the system forms a complete workflow from media upload to reviewed and anonymized output. The project shows how computer vision, redaction techniques, and user control can be integrated into a single workflow for detecting and anonymizing sensitive visual information.